\documentclass[11pt,a4paper]{article}

\usepackage[utf8]{inputenc}
\usepackage[vietnamese]{babel}
\usepackage[T5]{fontenc}
\usepackage[margin=1.7cm]{geometry}
\usepackage{amsmath,amssymb}
\usepackage{booktabs}
\usepackage{tabularx}
\usepackage{array}
\usepackage[table,xcdraw]{xcolor}
\usepackage{hyperref}
\usepackage{titlesec}
\usepackage{enumitem}
\usepackage{float}
\usepackage{tikz}
\usetikzlibrary{shapes.geometric, arrows.meta, positioning, trees, calc}

\hypersetup{
    colorlinks=true,
    linkcolor=blue!80!black,
    citecolor=blue!80!black,
    urlcolor=blue!80!black
}

\newcolumntype{Y}{>{\centering\arraybackslash}X}
\newcolumntype{L}{>{\raggedright\arraybackslash}X}

\begin{document}

% ==============================================================================
\section*{Bài tập 1: Website bán hàng}
% ==============================================================================

\subsection*{1. Yêu cầu đặc tả}
Website bán hàng có chính sách khuyến mãi như sau:
\begin{enumerate}[noitemsep]
    \item Khách mới giảm 10\%.
    \item Khách cũ giảm 15\%.
    \item Khách có thẻ giảm giá: 25\%.
\end{enumerate}

\subsection*{2. Xác định Conditions \& Actions}
\begin{itemize}[noitemsep]
    \item \textbf{Conditions (Điều kiện):}
    \begin{itemize}[noitemsep]
        \item \textbf{C1 (Loại khách hàng):} Khách mới / Khách cũ (hai trạng thái loại trừ).
        \item \textbf{C2 (Thẻ giảm giá):} Có (T) / Không (F).
    \end{itemize}
    \item \textbf{Actions (Hành động / Kết quả):}
    \begin{itemize}[noitemsep]
        \item \textbf{A1:} Giảm 10\%
        \item \textbf{A2:} Giảm 15\%
        \item \textbf{A3:} Giảm 25\%
    \end{itemize}
\end{itemize}

\textit{Quy tắc nghiệp vụ:} Thẻ giảm giá 25\% là mức ưu tiên áp dụng cao nhất, khi khách hàng có thẻ giảm giá thì được hưởng mức 25\% (áp dụng cho cả khách mới và khách cũ).

\subsection*{3. Hình vẽ: Sơ đồ Cây quyết định (Decision Tree)}
\begin{center}
\begin{tikzpicture}[
    level 1/.style={sibling distance=82mm, level distance=1.3cm},
    level 2/.style={sibling distance=34mm, level distance=1.4cm},
    decision/.style={rectangle, draw=blue!75, fill=blue!10, rounded corners=4pt, minimum height=2.1em, text centered, font=\bfseries\small, inner sep=4pt},
    outcome/.style={rectangle, draw=green!60!black, fill=green!15, rounded corners=4pt, minimum height=2.1em, text centered, font=\small\bfseries, inner sep=3pt, align=center},
    edge from parent/.style={draw=black!75, -Latex, thick},
    edge from parent path={(\tikzparentnode.south) -- ++(0,-0.3) -| (\tikzchildnode.north)}
]
\node [decision] {Loại khách hàng}
    child {
        node [decision] {Có thẻ giảm giá?}
        child {
            node [outcome] {Giảm 10\%\\(R1)}
            edge from parent node[left=2pt, font=\scriptsize\bfseries] {Không (F)}
        }
        child {
            node [outcome] {Giảm 25\%\\(R2)}
            edge from parent node[right=2pt, font=\scriptsize\bfseries] {Có (T)}
        }
        edge from parent node[above left=2pt, font=\scriptsize\bfseries] {Khách mới}
    }
    child {
        node [decision] {Có thẻ giảm giá?}
        child {
            node [outcome] {Giảm 15\%\\(R3)}
            edge from parent node[left=2pt, font=\scriptsize\bfseries] {Không (F)}
        }
        child {
            node [outcome] {Giảm 25\%\\(R4)}
            edge from parent node[right=2pt, font=\scriptsize\bfseries] {Có (T)}
        }
        edge from parent node[above right=2pt, font=\scriptsize\bfseries] {Khách cũ}
    };
\end{tikzpicture}
\end{center}

\subsection*{4. Bảng quyết định (Decision Table)}

\begin{table}[H]
\centering
\renewcommand{\arraystretch}{1.2}
\begin{tabularx}{\textwidth}{|l|l|Y|Y|Y|Y|}
\hline
\rowcolor{gray!20}
\multicolumn{2}{|c|}{\textbf{Thành phần}} & \textbf{R1} & \textbf{R2} & \textbf{R3} & \textbf{R4} \\ \hline
\textbf{Conditions} & \textbf{C1:} Loại khách hàng & Mới & Mới & Cũ & Cũ \\ \cline{2-6}
 & \textbf{C2:} Có thẻ giảm giá & F (Không) & T (Có) & F (Không) & T (Có) \\ \hline
\textbf{Actions} & \textbf{A1:} Giảm 10\% & \textbf{X} & & & \\ \cline{2-6}
 & \textbf{A2:} Giảm 15\% & & & \textbf{X} & \\ \cline{2-6}
 & \textbf{A3:} Giảm 25\% & & \textbf{X} & & \textbf{X} \\ \hline
\end{tabularx}
\caption*{Bảng quyết định cho bài toán khuyến mãi website}
\end{table}

\subsection*{5. Thiết kế ca kiểm thử (Test Cases)}

\begin{table}[H]
\centering
\renewcommand{\arraystretch}{1.2}
\begin{tabularx}{\textwidth}{|c|L|c|c|c|L|}
\hline
\rowcolor{blue!15}
\textbf{Test Case ID} & \textbf{Mô tả ca kiểm thử} & \textbf{Rule} & \textbf{Input: Loại khách} & \textbf{Input: Thẻ KM} & \textbf{Expected Output} \\ \hline
\textbf{TC01} & Khách mới không có thẻ giảm giá & R1 & Mới & Không & Giảm 10\% \\ \hline
\textbf{TC02} & Khách mới có thẻ giảm giá & R2 & Mới & Có & Giảm 25\% \\ \hline
\textbf{TC03} & Khách cũ không có thẻ giảm giá & R3 & Cũ & Không & Giảm 15\% \\ \hline
\textbf{TC04} & Khách cũ có thẻ giảm giá & R4 & Cũ & Có & Giảm 25\% \\ \hline
\end{tabularx}
\caption*{Danh sách ca kiểm thử Bài tập 1}
\end{table}

\newpage

% ==============================================================================
\section*{Bài 2: Bảo hiểm ô tô}
% ==============================================================================

\subsection*{1. Yêu cầu đặc tả}
Mức phí bảo hiểm ô tô được tính theo độ tuổi và giới tính:
\begin{enumerate}[noitemsep]
    \item Giới tính nữ, nhỏ hơn 65 tuổi: \$500.
    \item Giới tính nam, nhỏ hơn 25 tuổi: \$3000.
    \item Giới tính nam, từ 25 -- 64 tuổi: \$1000.
    \item Bất kỳ ai từ 65 tuổi trở lên: \$1500.
\end{enumerate}

\subsection*{2. Xác định Conditions \& Actions}
\begin{itemize}[noitemsep]
    \item \textbf{Conditions (Điều kiện):}
    \begin{itemize}[noitemsep]
        \item \textbf{C1 (Giới tính):} Nam / Nữ.
        \item \textbf{C2 (Độ tuổi):} $< 25$, $25 - 64$, $\ge 65$.
    \end{itemize}
    \item \textbf{Actions (Mức phí bảo hiểm):}
    \begin{itemize}[noitemsep]
        \item \textbf{A1:} Phí \$500
        \item \textbf{A2:} Phí \$1000
        \item \textbf{A3:} Phí \$1500
        \item \textbf{A4:} Phí \$3000
    \end{itemize}
\end{itemize}

\subsection*{3. Hình vẽ: Sơ đồ Cây quyết định (Decision Tree)}
\begin{center}
\begin{tikzpicture}[
    level 1/.style={sibling distance=86mm, level distance=1.3cm},
    decision/.style={rectangle, draw=blue!75, fill=blue!10, rounded corners=4pt, minimum height=2.1em, text centered, font=\bfseries\small, inner sep=4pt},
    outcome/.style={rectangle, draw=green!60!black, fill=green!15, rounded corners=4pt, minimum height=2.1em, text centered, font=\small\bfseries, inner sep=3pt, align=center},
    edge from parent/.style={draw=black!75, -Latex, thick},
    edge from parent path={(\tikzparentnode.south) -- ++(0,-0.3) -| (\tikzchildnode.north)}
]
\node [decision] {Giới tính}
    child {
        node [decision] {Độ tuổi}
        [sibling distance=34mm, level distance=1.4cm]
        child {
            node [outcome] {Phí \$500\\(R1, R2)}
            edge from parent node[left=2pt, font=\scriptsize\bfseries] {$< 65$}
        }
        child {
            node [outcome] {Phí \$1500\\(R3)}
            edge from parent node[right=2pt, font=\scriptsize\bfseries] {$\ge 65$}
        }
        edge from parent node[above left=2pt, font=\scriptsize\bfseries] {Nữ}
    }
    child {
        node [decision] {Độ tuổi}
        [sibling distance=26mm, level distance=1.4cm]
        child {
            node [outcome] {Phí \$3000\\(R4)}
            edge from parent node[left=2pt, font=\scriptsize\bfseries] {$< 25$}
        }
        child {
            node [outcome] {Phí \$1000\\(R5)}
            edge from parent node[fill=white, inner sep=2pt, font=\scriptsize\bfseries] {25 -- 64}
        }
        child {
            node [outcome] {Phí \$1500\\(R6)}
            edge from parent node[right=2pt, font=\scriptsize\bfseries] {$\ge 65$}
        }
        edge from parent node[above right=2pt, font=\scriptsize\bfseries] {Nam}
    };
\end{tikzpicture}
\end{center}

\subsection*{4. Bảng quyết định (Decision Table)}

\begin{table}[H]
\centering
\renewcommand{\arraystretch}{1.2}
\begin{tabularx}{\textwidth}{|l|l|Y|Y|Y|Y|Y|Y|}
\hline
\rowcolor{gray!20}
\multicolumn{2}{|c|}{\textbf{Thành phần}} & \textbf{R1} & \textbf{R2} & \textbf{R3} & \textbf{R4} & \textbf{R5} & \textbf{R6} \\ \hline
\textbf{Conditions} & \textbf{C1:} Giới tính & Nữ & Nữ & Nữ & Nam & Nam & Nam \\ \cline{2-8}
 & \textbf{C2:} Độ tuổi & $< 25$ & $25 - 64$ & $\ge 65$ & $< 25$ & $25 - 64$ & $\ge 65$ \\ \hline
\textbf{Actions} & \textbf{A1:} Phí \$500 & \textbf{X} & \textbf{X} & & & & \\ \cline{2-8}
 & \textbf{A2:} Phí \$1000 & & & & & \textbf{X} & \\ \cline{2-8}
 & \textbf{A3:} Phí \$1500 & & & \textbf{X} & & & \textbf{X} \\ \cline{2-8}
 & \textbf{A4:} Phí \$3000 & & & & \textbf{X} & & \\ \hline
\end{tabularx}
\caption*{Bảng quyết định cho bài toán phí bảo hiểm ô tô}
\end{table}

\subsection*{5. Thiết kế ca kiểm thử (Test Cases)}

\begin{table}[H]
\centering
\renewcommand{\arraystretch}{1.2}
\begin{tabularx}{\textwidth}{|c|L|c|c|c|L|}
\hline
\rowcolor{blue!15}
\textbf{Test Case ID} & \textbf{Mô tả ca kiểm thử} & \textbf{Rule} & \textbf{Input: Giới tính} & \textbf{Input: Tuổi} & \textbf{Expected Output} \\ \hline
\textbf{TC01} & Nữ, tuổi dưới 25 & R1 & Nữ & 20 & Mức phí: \$500 \\ \hline
\textbf{TC02} & Nữ, tuổi từ 25 đến 64 & R2 & Nữ & 30 & Mức phí: \$500 \\ \hline
\textbf{TC03} & Nữ, từ 65 tuổi trở lên & R3 & Nữ & 68 & Mức phí: \$1500 \\ \hline
\textbf{TC04} & Nam, tuổi dưới 25 & R4 & Nam & 22 & Mức phí: \$3000 \\ \hline
\textbf{TC05} & Nam, tuổi từ 25 đến 64 & R5 & Nam & 40 & Mức phí: \$1000 \\ \hline
\textbf{TC06} & Nam, từ 65 tuổi trở lên & R6 & Nam & 70 & Mức phí: \$1500 \\ \hline
\end{tabularx}
\caption*{Danh sách ca kiểm thử Bài tập 2}
\end{table}

\newpage

% ==============================================================================
\section*{Bài 3: Kiểm tra điều kiện dự thi của sinh viên}
% ==============================================================================

\subsection*{1. Yêu cầu đặc tả}
Mô đun kiểm tra điều kiện dự thi của sinh viên gồm:
\begin{enumerate}[noitemsep]
    \item Nếu sinh viên đi học $\ge 80\%$ số buổi, điểm giữa kì $\ge 1$, điểm bài tập lớn $\ge 1$ thì được dự thi.
    \item Nếu sinh viên đủ điều kiện dự thi và có điểm thi giữa kì $= 10$ và điểm bài tập lớn $= 10$ thì được miễn thi.
    \item Trường hợp không thỏa mãn các điều kiện dự thi thì không được dự thi (cấm thi).
\end{enumerate}

\subsection*{2. Xác định Conditions \& Actions}
\begin{itemize}[noitemsep]
    \item \textbf{Conditions (Điều kiện):}
    \begin{itemize}[noitemsep]
        \item \textbf{C1 (Tỷ lệ đi học):} $\ge 80\%$ (T: Đúng, F: Sai).
        \item \textbf{C2 (Điểm giữa kì - GK):} $< 1$ (điểm liệt), $1 \le \text{GK} < 10$, $\text{GK} = 10$.
        \item \textbf{C3 (Điểm bài tập lớn - BTL):} $< 1$ (điểm liệt), $1 \le \text{BTL} < 10$, $\text{BTL} = 10$.
    \end{itemize}
    \item \textbf{Actions (Kết quả xét thi):}
    \begin{itemize}[noitemsep]
        \item \textbf{A1:} Không được dự thi (Cấm thi)
        \item \textbf{A2:} Được dự thi
        \item \textbf{A3:} Được miễn thi
    \end{itemize}
\end{itemize}

\subsection*{3. Hình vẽ: Sơ đồ Luồng quyết định (Decision Flowchart)}
\begin{center}
\begin{tikzpicture}[
    auto,
    startstop/.style={rectangle, draw=purple!80, fill=purple!15, rounded corners=6pt, text width=6.5em, text centered, minimum height=2.2em, font=\bfseries\small},
    decision/.style={diamond, aspect=1.8, draw=blue!75, fill=blue!12, text width=7.0em, text badly centered, inner sep=1pt, font=\bfseries\scriptsize},
    block_fail/.style={rectangle, draw=red!75, fill=red!15, rounded corners=4pt, text width=9.2em, text centered, minimum height=2.3em, font=\bfseries\scriptsize},
    block_pass/.style={rectangle, draw=green!60!black, fill=green!15, rounded corners=4pt, text width=9.2em, text centered, minimum height=2.3em, font=\bfseries\scriptsize},
    block_exempt/.style={rectangle, draw=orange!80!black, fill=yellow!20, rounded corners=4pt, text width=9.2em, text centered, minimum height=2.3em, font=\bfseries\scriptsize},
    line/.style={draw=black!75, -Latex, thick}
]
    % Nodes with clear spacing to avoid overlapping
    \node [startstop] (start) at (0, 0) {Bắt đầu};
    \node [decision] (d1) at (0, -1.8) {Đi học $\ge 80\%$?};
    \node [decision] (d2) at (0, -4.3) {Điểm GK $\ge 1$ và\\Điểm BTL $\ge 1$?};
    \node [decision] (d3) at (0, -6.8) {Điểm GK $= 10$ và\\Điểm BTL $= 10$?};
    
    \node [block_fail] (r1) at (5.5, -1.8) {Không được dự thi\\(R1)};
    \node [block_fail] (r23) at (5.5, -4.3) {Không được dự thi\\(R2, R3)};
    \node [block_pass] (r45) at (5.5, -6.8) {Được dự thi\\(R4, R5)};
    \node [block_exempt] (r6) at (0, -8.8) {Được miễn thi (R6)};

    % Lines
    \path [line] (start) -- (d1);
    \path [line] (d1) -- node[right=4pt, font=\scriptsize\bfseries] {Đúng (T)} (d2);
    \path [line] (d1) -- node[above=2pt, font=\scriptsize\bfseries] {Sai (F)} (r1);
    
    \path [line] (d2) -- node[right=4pt, font=\scriptsize\bfseries] {Đúng (T)} (d3);
    \path [line] (d2) -- node[above=2pt, font=\scriptsize\bfseries] {Sai (F)} (r23);
    
    \path [line] (d3) -- node[above=2pt, font=\scriptsize\bfseries] {Sai (F)} (r45);
    \path [line] (d3) -- node[right=4pt, font=\scriptsize\bfseries] {Đúng (T)} (r6);
\end{tikzpicture}
\end{center}

\subsection*{4. Bảng quyết định (Decision Table)}

\begin{table}[H]
\centering
\renewcommand{\arraystretch}{1.25}
\begin{tabularx}{\textwidth}{|l|l|Y|Y|Y|Y|Y|Y|}
\hline
\rowcolor{gray!20}
\multicolumn{2}{|c|}{\textbf{Thành phần}} & \textbf{R1} & \textbf{R2} & \textbf{R3} & \textbf{R4} & \textbf{R5} & \textbf{R6} \\ \hline
\textbf{Conditions} & \textbf{C1:} Đi học $\ge 80\%$ & \textbf{F} & \textbf{T} & \textbf{T} & \textbf{T} & \textbf{T} & \textbf{T} \\ \cline{2-8}
 & \textbf{C2:} Điểm giữa kì (GK) & -- & $< 1$ & $\ge 1$ & $1 \le \text{GK} < 10$ & $10$ & $10$ \\ \cline{2-8}
 & \textbf{C3:} Điểm BTL & -- & -- & $< 1$ & $1 \le \text{BTL} \le 10$ & $1 \le \text{BTL} < 10$ & $10$ \\ \hline
\textbf{Actions} & \textbf{A1:} Không được dự thi & \textbf{X} & \textbf{X} & \textbf{X} & & & \\ \cline{2-8}
 & \textbf{A2:} Được dự thi & & & & \textbf{X} & \textbf{X} & \\ \cline{2-8}
 & \textbf{A3:} Được miễn thi & & & & & & \textbf{X} \\ \hline
\end{tabularx}
\caption*{Bảng quyết định cho quy tắc điều kiện dự thi của sinh viên}
\end{table}

\noindent\textit{Diễn giải các quy tắc:}
\begin{itemize}[noitemsep]
    \item \textbf{R1:} Đi học $< 80\%$ $\rightarrow$ Bị cấm thi (bất kể điểm số).
    \item \textbf{R2:} Đi học đủ nhưng Điểm giữa kì $< 1$ $\rightarrow$ Bị cấm thi.
    \item \textbf{R3:} Đi học đủ, GK $\ge 1$ nhưng Điểm BTL $< 1$ $\rightarrow$ Bị cấm thi.
    \item \textbf{R4:} Đi học đủ, cả hai điểm $\ge 1$ (trong đó $\text{GK} < 10$) $\rightarrow$ Được dự thi.
    \item \textbf{R5:} Đi học đủ, $\text{GK} = 10$ nhưng $\text{BTL} < 10$ ($\ge 1$) $\rightarrow$ Được dự thi.
    \item \textbf{R6:} Đi học đủ, đồng thời $\text{GK} = 10$ và $\text{BTL} = 10$ $\rightarrow$ Được miễn thi.
\end{itemize}

\subsection*{5. Thiết kế ca kiểm thử (Test Cases)}

\begin{table}[H]
\centering
\renewcommand{\arraystretch}{1.2}
\begin{tabularx}{\textwidth}{|c|L|c|c|c|c|L|}
\hline
\rowcolor{blue!15}
\textbf{TC ID} & \textbf{Mô tả ca kiểm thử} & \textbf{Rule} & \textbf{Tỷ lệ học} & \textbf{Điểm GK} & \textbf{Điểm BTL} & \textbf{Expected Output} \\ \hline
\textbf{TC01} & Nghỉ quá số buổi cho phép & R1 & 75\% & 8.0 & 9.0 & Không được dự thi \\ \hline
\textbf{TC02} & Đi học đủ, GK bị điểm liệt ($< 1$) & R2 & 90\% & 0.5 & 8.0 & Không được dự thi \\ \hline
\textbf{TC03} & Đi học đủ, BTL bị điểm liệt ($< 1$) & R3 & 85\% & 7.0 & 0.0 & Không được dự thi \\ \hline
\textbf{TC04} & Đủ điều kiện dự thi thông thường & R4 & 80\% & 6.5 & 8.0 & Được dự thi \\ \hline
\textbf{TC05} & Đủ điều kiện, GK = 10 nhưng BTL < 10 & R5 & 95\% & 10.0 & 8.5 & Được dự thi \\ \hline
\textbf{TC06} & Đạt điểm tối đa cả hai môn & R6 & 100\% & 10.0 & 10.0 & Được miễn thi \\ \hline
\end{tabularx}
\caption*{Danh sách ca kiểm thử Bài tập 3}
\end{table}

\newpage

% ==============================================================================
\section*{Bài tập: Kiểm thử dòng điều khiển (CFG) \& Đường cơ sở}
% ==============================================================================

\subsection*{1. Yêu cầu đặc tả \& Đoạn mã nguồn}
\noindent\textbf{Yêu cầu:}
\begin{itemize}[noitemsep]
    \item Vẽ sơ đồ dòng điều khiển (Control Flow Graph -- CFG) cho đoạn mã.
    \item Xác định các đường đi độc lập theo phương pháp kiểm thử đường cơ sở (Basis Path Testing).
    \item Thiết kế bộ test đảm bảo bao phủ tất cả các đường đi độc lập.
\end{itemize}

\noindent\textbf{Đoạn mã nguồn phân loại tam giác (quan hệ giữa $a, b, c$):}
\begin{center}
\begin{minipage}{0.9\textwidth}
\begin{tabular}{|r l|}
\hline
\rowcolor{gray!15}
\textbf{Dòng} & \textbf{Mã lệnh (Source Code)} \\ \hline
1: & \texttt{String classifyTriangle(int a, int b, int c) \{} \\
2: & \quad \texttt{String type;} \\
3: & \quad \texttt{if (a + b <= c || a + c <= b || b + c <= a) \{} \\
4: & \qquad \texttt{type = "Không phải tam giác";} \\
5: & \quad \texttt{\} else if (a == b \&\& b == c) \{} \\
6: & \qquad \texttt{type = "Tam giác đều";} \\
7: & \quad \texttt{\} else if (a == b || b == c || a == c) \{} \\
8: & \qquad \texttt{type = "Tam giác cân";} \\
9: & \quad \texttt{\} else \{} \\
10: & \qquad \texttt{type = "Tam giác thường";} \\
11: & \quad \texttt{\}} \\
12: & \quad \texttt{return type;} \\
13: & \texttt{\}} \\ \hline
\end{tabular}
\end{minipage}
\end{center}

\subsection*{2. Danh sách các nút (Nodes) và cung (Edges) của đồ thị CFG}
\begin{itemize}[noitemsep]
    \item \textbf{Nút (Nodes):}
    \begin{itemize}[noitemsep]
        \item \textbf{Nút (1, 2):} Bắt đầu hàm và khai báo biến \texttt{type}.
        \item \textbf{Nút (3):} Nút quyết định 1 -- Kiểm tra BĐT tam giác ($a + b \le c \lor a + c \le b \lor b + c \le a$).
        \item \textbf{Nút (4):} Xử lý gán \texttt{type = "Không phải tam giác"}.
        \item \textbf{Nút (5):} Nút quyết định 2 -- Kiểm tra 3 cạnh bằng nhau ($a == b \land b == c$).
        \item \textbf{Nút (6):} Xử lý gán \texttt{type = "Tam giác đều"}.
        \item \textbf{Nút (7):} Nút quyết định 3 -- Kiểm tra có 2 cạnh bằng nhau ($a == b \lor b == c \lor a == c$).
        \item \textbf{Nút (8):} Xử lý gán \texttt{type = "Tam giác cân"}.
        \item \textbf{Nút (10):} Xử lý gán \texttt{type = "Tam giác thường"}.
        \item \textbf{Nút (12, 13):} Lệnh \texttt{return type} và kết thúc hàm.
    \end{itemize}
    \item \textbf{Tổng số nút:} $N = 8$ (coi các nút tuần tự liền kề là 1 nút: $\{1,2\}$, $\{3\}$, $\{4\}$, $\{5\}$, $\{6\}$, $\{7\}$, $\{8\}$, $\{10\}$, $\{12,13\}$).
    \item \textbf{Tổng số cung:} $E = 10$ cung có hướng.
\end{itemize}

\subsection*{3. Hình vẽ: Sơ đồ dòng điều khiển (Control Flow Graph -- CFG)}
\begin{center}
\begin{tikzpicture}[
    auto,
    cfgnode/.style={circle, draw=blue!80, fill=blue!10, thick, minimum size=0.85cm, font=\bfseries\small},
    decnode/.style={circle, draw=orange!90!black, fill=orange!20, thick, minimum size=0.85cm, font=\bfseries\small},
    endnode/.style={circle, double, draw=purple!80, fill=purple!15, thick, minimum size=0.95cm, font=\bfseries\small},
    line/.style={draw=black!75, -Latex, thick},
    region/.style={draw=blue!50, fill=white, rounded corners=2pt, font=\bfseries\footnotesize\color{blue!80!black}, inner sep=2pt}
]
    % Central spine
    \node [cfgnode] (n12) at (0, 0) {1,2};
    \node [decnode] (n3) at (0, -1.4) {3};
    \node [decnode] (n5) at (0, -3.0) {5};
    \node [decnode] (n7) at (0, -4.6) {7};
    \node [cfgnode] (n10) at (0, -6.2) {10};
    \node [endnode] (n1213) at (0, -7.8) {12,13};

    % Right branches
    \node [cfgnode] (n4) at (3.5, -2.1) {4};
    \node [cfgnode] (n6) at (3.5, -3.7) {6};
    \node [cfgnode] (n8) at (3.5, -5.3) {8};

    % Edges
    \path [line] (n12) -- (n3);
    \path [line] (n3) -- node[above=2pt, font=\footnotesize\bfseries] {True} (n4);
    \path [line] (n3) -- node[left=2pt, font=\footnotesize\bfseries] {False} (n5);
    
    \path [line] (n5) -- node[above=2pt, font=\footnotesize\bfseries] {True} (n6);
    \path [line] (n5) -- node[left=2pt, font=\footnotesize\bfseries] {False} (n7);
    
    \path [line] (n7) -- node[above=2pt, font=\footnotesize\bfseries] {True} (n8);
    \path [line] (n7) -- node[left=2pt, font=\footnotesize\bfseries] {False} (n10);
    
    \path [line] (n10) -- (n1213);

    % Lines to exit with different routing offsets and distinct arrowheads
    \draw [line] (n4.east) -- ++(2.0, 0) |- ($(n1213.east) + (0, 0.35)$);
    \draw [line] (n6.east) -- ++(1.4, 0) |- (n1213.east);
    \draw [line] (n8.east) -- ++(0.8, 0) |- ($(n1213.east) + (0, -0.35)$);

    % Region annotations placed in open areas with white backgrounds
    \node [region] at (1.6, -2.3) {$R_1$};
    \node [region] at (1.6, -3.9) {$R_2$};
    \node [region] at (1.6, -5.5) {$R_3$};
    \node [region] at (-2.0, -0.8) {$R_4$ (Miền ngoài)};
\end{tikzpicture}
\end{center}

\subsection*{4. Tính độ phức tạp Cyclomatic $V(G)$}
Độ phức tạp Cyclomatic $V(G)$ được tính bằng 3 công thức tương đương:
\begin{enumerate}[noitemsep]
    \item \textbf{Theo số cạnh và số nút:}
    \[
    V(G) = E - N + 2P = 10 - 8 + 2(1) = 4
    \]
    (với $E=10$ cạnh, $N=8$ nút, $P=1$ thành phần liên thông).
    \item \textbf{Theo số nút điều kiện (Predicate Nodes):}
    Đồ thị có 3 nút rẽ nhánh điều kiện là $\{3, 5, 7\}$, do đó $P_{red} = 3$:
    \[
    V(G) = P_{red} + 1 = 3 + 1 = 4
    \]
    \item \textbf{Theo số vùng/miền độc lập (Regions):}
    Đồ thị có 3 miền đóng kín ($R_1, R_2, R_3$) và 1 miền mở vô hạn bên ngoài ($R_4$):
    \[
    V(G) = \text{Số miền} = 4
    \]
\end{enumerate}
$\Rightarrow$ \textbf{Kết luận:} Số đường đi độc lập cơ sở tối thiểu cần thiết để bao phủ là \textbf{4 đường đi}.

\subsection*{5. Xác định các đường đi độc lập cơ sở (Basis Paths)}
Từ đồ thị dòng điều khiển, tập các đường đi độc lập cơ sở bao gồm:
\begin{itemize}[noitemsep]
    \item \textbf{Path 1:} $(1,2) \to 3 \to 4 \to (12,13)$ \\
    \textit{Ý nghĩa:} Không thỏa mãn bất đẳng thức tam giác $\rightarrow$ Không phải tam giác.
    \item \textbf{Path 2:} $(1,2) \to 3 \to 5 \to 6 \to (12,13)$ \\
    \textit{Ý nghĩa:} Là tam giác và cả 3 cạnh bằng nhau $\rightarrow$ Tam giác đều.
    \item \textbf{Path 3:} $(1,2) \to 3 \to 5 \to 7 \to 8 \to (12,13)$ \\
    \textit{Ý nghĩa:} Là tam giác, không đều nhưng có 2 cạnh bằng nhau $\rightarrow$ Tam giác cân.
    \item \textbf{Path 4:} $(1,2) \to 3 \to 5 \to 7 \to 10 \to (12,13)$ \\
    \textit{Ý nghĩa:} Là tam giác, không đều và không cân $\rightarrow$ Tam giác thường.
\end{itemize}

\subsection*{6. Thiết kế bộ ca kiểm thử (Basis Path Test Cases)}

Bảng ca kiểm thử đảm bảo bao phủ 100\% các đường đi độc lập cơ sở:

\begin{table}[H]
\centering
\renewcommand{\arraystretch}{1.3}
\begin{tabularx}{\textwidth}{|c|L|c|c|c|L|L|}
\hline
\rowcolor{blue!15}
\textbf{TC ID} & \textbf{Đường đi bao phủ} & $\mathbf{a}$ & $\mathbf{b}$ & $\mathbf{c}$ & \textbf{Tổ hợp điều kiện} & \textbf{Expected Output} \\ \hline
\textbf{TC01} & \textbf{Path 1:} $(1,2) \to 3 \to 4 \to 12$ & 1 & 2 & 5 & $1 + 2 \le 5$ (Sai BĐT tam giác) & Không phải tam giác \\ \hline
\textbf{TC02} & \textbf{Path 2:} $(1,2) \to 3 \to 5 \to 6 \to 12$ & 4 & 4 & 4 & $a = b = c = 4$ & Tam giác đều \\ \hline
\textbf{TC03} & \textbf{Path 3:} $(1,2) \to 3 \to 5 \to 7 \to 8 \to 12$ & 3 & 3 & 5 & $a = b = 3 \ne c$ & Tam giác cân \\ \hline
\textbf{TC04} & \textbf{Path 4:} $(1,2) \to 3 \to 5 \to 7 \to 10 \to 12$ & 3 & 4 & 5 & Thỏa BĐT tam giác, $a \ne b \ne c$ & Tam giác thường \\ \hline
\end{tabularx}
\caption*{Bộ ca kiểm thử bao phủ tất cả các đường đi độc lập}
\end{table}

\end{document}
